\section{Related Work}

\noindent\textbf{Tactile Sensing Technologies.} Tactile sensors are capable of direct, contact-level feedback on force, slip, and material properties~\cite{dahiya2009tactile, li2020tactile}. 
Recently, a diverse set of transduction principles has been extensively explored.
Vision-based tactile sensors, e.g., GelSight~\cite{yuan2017gelsight}, DIGIT~\cite{lambeta2020digit}, and TacTip~\cite{ward2018tactip}, capture elastomer deformation with an embedded camera and extract high-resolution geometry and force fields through computer-vision pipelines.
However, they are typically constrained by camera frame rate, image-processing overhead, and bulky optical packaging.
Capacitive tactile sensors measure deformation-induced capacitance changes~\cite{schmitz2011methods}.
However, this approach falls short in sensing fine-grained force information and is susceptible to parasitic capacitance, humidity, and temperature drift.
Piezoresistive tactile sensors exploit force-induced resistance changes~\cite{stassi2014flexible}, thus having lower costs.
However, this approach suffers from hysteresis, drift, and limited long-term stability.
In contrast, magnetic tactile sensors transduce contact forces by tracking the displacement of an embedded magnet with a Hall-effect or magnetoresistive array~\cite{bhirangi2021reskin, yan2021softmagnetic}, enabling fine-grained three-dimensional force decoupling, compact form factors, robustness to optical occlusion, and low fabrication cost.

\noindent\textbf{Physical-Layer Signal Injection Attacks.} 
A broad class of prior work has demonstrated that adversaries can manipulate embedded sensors by injecting crafted physical signals, bypassing higher-level software defenses.
Kune et al.~\cite{kune2013ghost} first systematically characterized electromagnetic signal injection against analog sensors, showing that EMI can induce spurious readings in implantable cardiac devices and consumer microphones without any physical contact.
Trippel et al.~\cite{trippel2017walnut} demonstrated that MEMS accelerometers in wearables and drones can be driven to arbitrary output values via acoustic resonance injection, corrupting the inertial measurements that autonomous systems depend on.
Ji et al.~\cite{ji2021poltergeist} extended this idea to vision pipelines, using acoustic signals to cause resonance-induced blur in camera image sensors, thereby degrading the perception of computer-vision-based autonomous vehicles.
Xu et al.~\cite{xu2018analyzing} analyzed signal injection vulnerabilities in ultrasonic sensors used for obstacle detection in autonomous vehicles, demonstrating that spoofed echo signals can suppress or fabricate obstacles.
Zhang et al.~\cite{zhang2017dolphinattack} and Yan et al.~\cite{yan2019feasibility} showed that voice assistants are vulnerable to inaudible ultrasonic command injection, which is imperceptible to humans yet recognized as valid speech by microphone front-ends.
{\name} shares this adversarial philosophy but targets a qualitatively different modality---magnetic tactile sensing---and introduces a uniquely passive, \emph{deploy-and-forget} attack vector using battery-free permanent magnet tags, whereas prior injection attacks typically require active electronic signal generators in close proximity.

\noindent\textbf{Security of Robotic Perception and Tactile Sensing.}
As embodied AI systems are expanding to safety-critical settings, the security of robotic perception has attracted increasing attention. 
Prior work has examined adversarial attacks on visual pipelines~\cite{ji2021poltergeist}, ultrasonic obstacle sensors~\cite{xu2018analyzing}, and inertial measurement units~\cite{trippel2017walnut}.
However, limited efforts have been made into the potential security issues of tactile sensing.
Existing anti-interference mechanisms in magnetic tactile sensors, such as gradiometric differential arrays~\cite{le2022gradiometric}, are only capable of overcoming common interfering magnetic fields. 
To the best of our knowledge, {\name} is the first effort to characterize and exploit the security loophole in magnetic tactile sensors. 
{\name} demonstrates the end-to-end hijacking performance of robotic closed-loop force control with an attack success rate exceeding 95\%.
